%  00_abstract.tex
\begin{abstract}
We present an open-set individual face recognition system for fruit bats, designed to
support longitudinal behavioural studies where every animal must be identified without
prior appearance in the training set.
Our primary methodological contribution is a rigorous \emph{identity-disjoint} evaluation
protocol: train, validation, and test splits are partitioned at the level of individual
identities so that no bat ever appears in more than one split.
We show that the widely-used alternative---random pair selection---inflates performance
dramatically: a Siamese network that achieves $F_1 = 0.99$ under the random-pair protocol
drops to $F_1 = 0.74$ under identity-disjoint evaluation on the same images.

We compare three deep-learning model families---a pair-based Siamese convolutional
network and two embedding-based angular-margin architectures (ArcFace and AdaFace built
on a ResNet-50 backbone)---across a factorial design covering two fruit-bat species
(\rousettus and \mauritius), three background conditions (green screen, original, and
random compositing), and controlled identity counts.
On the rousettus dataset (12 individuals, 1\,059 images), the Siamese network achieves a
test \rocauc of $0.809$ versus $0.541$ and $0.507$ for ArcFace and AdaFace respectively.
On mauritius (11 individuals, 1\,119 images) the ordering is preserved, with the Siamese
network reaching $0.803$--$0.925$ depending on background condition.
We trace the embedding-model shortfall to the \emph{small-identity generalisation wall}:
with fewer than twelve training identities the angular separation learned during closed-set
training does not transfer to unseen test identities, even when training accuracy
approaches 100\%.

All experiments are fully reproducible via a single-command pipeline and are tracked in
MLflow.
\end{abstract}
